% Document compilable avec: pdflatex scRAW_Default_Configuration.tex
% =============================================================================
\documentclass[11pt,a4paper]{article}

% --- Packages ----------------------------------------------------------------
\usepackage[utf8]{inputenc}
\usepackage[T1]{fontenc}
\usepackage[french]{babel}
\usepackage{lmodern}
\usepackage{amsmath, amssymb}
\usepackage{graphicx}
\usepackage{geometry}
\geometry{margin=2.2cm}
\usepackage{hyperref}
\usepackage{tabularx}
\usepackage{booktabs}
\usepackage{longtable}
\usepackage{multirow}
\usepackage{array}
\usepackage{tikz}
\usetikzlibrary{shapes.geometric, arrows, arrows.meta, positioning, fit, calc, backgrounds, decorations.pathreplacing, patterns, matrix}
\usepackage{pgfplots}
\pgfplotsset{compat=1.18}
\usepackage{tcolorbox}
\tcbuselibrary{skins, breakable}
\usepackage{caption}
\usepackage{float}
\usepackage{enumitem}
\usepackage{xcolor}

% --- Custom colors -----------------------------------------------------------
\definecolor{noteblue}{HTML}{0969DA}
\definecolor{notebg}{HTML}{DDF4FF}
\definecolor{tipgreen}{HTML}{1A7F37}
\definecolor{tipbg}{HTML}{D3F5D7}
\definecolor{importantorange}{HTML}{9A6700}
\definecolor{importantbg}{HTML}{FFF8C5}
\definecolor{cautionred}{HTML}{CF222E}
\definecolor{cautionbg}{HTML}{FFD8D3}
\definecolor{codegray}{HTML}{F6F8FA}
\definecolor{latentcolor}{HTML}{FF9800}
\definecolor{phaseblue}{HTML}{1565C0}
\definecolor{phaseorange}{HTML}{E65100}
\definecolor{phasegreen}{HTML}{2E7D32}
\definecolor{failred}{HTML}{D32F2F}
\definecolor{successgreen}{HTML}{388E3C}
\definecolor{dannred}{HTML}{C62828}
\definecolor{weightpurple}{HTML}{6A1B9A}
\definecolor{tripletyellow}{HTML}{F9A825}
\definecolor{nbblue}{HTML}{1976D2}
\definecolor{enccolor}{HTML}{1E88E5}
\definecolor{deccolor}{HTML}{E53935}
\definecolor{clustergreen}{HTML}{43A047}

% --- Custom boxes ------------------------------------------------------------
\newtcolorbox{notebox}{
  colback=notebg, colframe=noteblue, fonttitle=\bfseries,
  title={\textcolor{noteblue}{$\gg$ Note}}, breakable, left=4mm
}
\newtcolorbox{tipbox}{
  colback=tipbg, colframe=tipgreen, fonttitle=\bfseries,
  title={\textcolor{tipgreen}{$\gg$ Astuce}}, breakable, left=4mm
}
\newtcolorbox{importantbox}{
  colback=importantbg, colframe=importantorange, fonttitle=\bfseries,
  title={\textcolor{importantorange}{$\gg$ Important}}, breakable, left=4mm
}
\newtcolorbox{cautionbox}{
  colback=cautionbg, colframe=cautionred, fonttitle=\bfseries,
  title={\textcolor{cautionred}{$\gg$ Attention}}, breakable, left=4mm
}
\newtcolorbox{parambox}[1]{
  colback=codegray, colframe=gray!50, fonttitle=\bfseries\ttfamily,
  title={#1}, breakable, left=4mm
}

% --- Hyperref setup ----------------------------------------------------------
\hypersetup{
  colorlinks=true,
  linkcolor=noteblue,
  urlcolor=noteblue,
  citecolor=noteblue,
}

% --- Title -------------------------------------------------------------------
\title{\textbf{scRAW --- Document Explicatif}\\[4pt]
{\Large Fonctionnement de l'autoencoder et du weighting system}}
\author{\'Equipe SCRBenchmark}
\date{Mars 2026}

\begin{document}

\maketitle

\begin{abstract}
Ce document d\'ecrit le fonctionnement de \textbf{scRAW} (Single-Cell Rare-population Autoencoder Weighting), un autoencoder d\'eterministe pour le clustering de donn\'ees scRNA-seq avec une attention particuli\`ere aux populations cellulaires rares. Il pr\'esente les \'etapes du pipeline : preprocessing, architecture r\'eseau, syst\`eme de pond\'eration dynamique bicrit\`ere et r\'egularisation par Triplet Loss. Le module optionnel de correction de batch (DANN) est \'egalement d\'ecrit.
\end{abstract}

\vspace{0.5cm}
\tableofcontents
\newpage

% =============================================================================
\section{Description du pipeline}
\label{sec:pipeline}
% =============================================================================

scRAW est un \textbf{autoencoder d\'eterministe} pour le clustering de donn\'ees scRNA-seq. Il n'utilise pas de sampling dans le latent space : chaque cellule produit un vecteur $z$ de 128 dimensions.

\begin{figure}[H]
\centering
\resizebox{\linewidth}{!}{%
\begin{tikzpicture}[
  node distance=0.9cm,
  stage/.style={rectangle, draw, thick, rounded corners=8pt, align=center,
    minimum width=2.8cm, minimum height=1.5cm, font=\small},
  minitag/.style={rectangle, rounded corners=3pt, fill=#1, font=\tiny\bfseries\color{white},
    inner sep=2pt, minimum height=0.4cm},
  arr/.style={-{Stealth[length=3mm]}, thick, gray!70},
  darr/.style={-{Stealth[length=3mm]}, thick, dashed, gray!50},
]

% Main pipeline stages
\node[stage, fill=gray!8] (data) {\textbf{Raw counts}\\$n \times G$ genes};

\node[stage, fill=blue!6, right=0.9cm of data] (preproc) {\textbf{Preprocessing}\\Norm + HVG\\+ Scale};

\node[stage, fill=phaseblue!10, right=0.9cm of preproc] (warmup) {\textbf{Phase 1}\\Warm-up\\30 ep.};

\node[stage, fill=phaseorange!12, right=0.9cm of warmup] (weighted) {\textbf{Phase 2}\\Weighted\\+ Triplet\\(30--120)};

\node[stage, fill=clustergreen!12, right=0.9cm of weighted] (cluster) {\textbf{Phase 3}\\HDBSCAN\\c4s2, EOM};

\node[stage, fill=successgreen!15, right=0.9cm of cluster] (out) {$k$ clusters\\output};

% Tags
\node[minitag=gray!60, above=0.15cm of data.north] {INPUT};
\node[minitag=phaseblue, above=0.15cm of warmup.north] {UNIFORM NB (0--29)};
\node[minitag=phaseorange, above=0.15cm of weighted.north] {WEIGHTED NB};
\node[minitag=successgreen, above=0.15cm of out.north] {OUTPUT};

% Arrows
\draw[arr] (data) -- node[above, font=\tiny]{\texttt{log1p}} (preproc);
\draw[arr] (preproc) -- node[below, font=\tiny]{Encoder input : log1p} (warmup);
\draw[arr] (warmup) -- node[above, font=\tiny]{Weights} node[below, font=\tiny]{Leiden} (weighted);
\draw[arr] (weighted) -- node[above, font=\tiny]{$z_{128}$} (cluster);
\draw[arr, clustergreen] (cluster) -- (out);

% (DANN optionnel, activable en multi-batch)

% Weight update annotation (internal to phase 2, every 10 epochs)
\node[font=\tiny\bfseries, text=weightpurple, below=0.18cm of weighted.south, align=center]
  {\textit{Pseudo-labels Leiden} $\to$ $w_i$ ($\Delta$10 ép., EMA)};

\end{tikzpicture}%
}
\caption{Pipeline global scRAW. Preprocessing Scanpy (log1p), training en deux phases (warm-up NB uniforme, puis weighted phase avec Triplet Loss), puis clustering HDBSCAN sur le latent space. Le module DANN de correction de batch peut \^etre activ\'e en configuration multi-batch.}
\label{fig:pipeline_global}
\end{figure}


% =============================================================================
\section{Network architecture}
\label{sec:architecture}
% =============================================================================

\subsection{Structure de l'autoencoder}

Autoencoder sym\'etrique avec BatchNorm, LeakyReLU et Dropout. Le latent space est une pure linear projection (128D).

\begin{figure}[H]
\centering
\resizebox{\linewidth}{!}{%
\begin{tikzpicture}[
  layer/.style={rectangle, draw, thick, rounded corners=4pt, minimum width=1.8cm,
    minimum height=0.9cm, align=center, font=\footnotesize},
  op/.style={rectangle, rounded corners=2pt, fill=#1, font=\tiny\bfseries\color{white},
    inner sep=1.5pt, minimum height=0.35cm},
  arr/.style={-{Stealth[length=2.5mm]}, thick},
  brace/.style={decorate, decoration={brace, amplitude=6pt, raise=2pt}},
]

% ---- Encoder ----
\node[layer, fill=enccolor!10, text width=1.6cm] (input) {\textbf{Input}\\$n_{genes}$};

\node[layer, fill=enccolor!15, right=1.8cm of input, text width=1.6cm] (e1) {\textbf{512}};
\node[op=enccolor, above=0.08cm of e1.north] {Lin + BN + LReLU + Drop};

\node[layer, fill=enccolor!20, right=1.5cm of e1, text width=1.6cm] (e2) {\textbf{256}};
\node[op=enccolor, above=0.08cm of e2.north] {Lin + BN + LReLU + Drop};

\node[layer, fill=enccolor!25, right=1.5cm of e2, text width=1.6cm] (e3) {\textbf{128}};
\node[op=enccolor, above=0.08cm of e3.north] {Lin + BN + LReLU + Drop};

% ---- Latent ----
\node[layer, fill=latentcolor!25, right=1.5cm of e3, text width=1.8cm, minimum height=1.2cm,
  line width=1.5pt, draw=latentcolor] (z) {\textbf{Latent space}\\$z = 128$\\(Pure linear)};

% ---- Decoder ----
\node[layer, fill=deccolor!10, right=1.5cm of z, text width=1.6cm] (d1) {\textbf{128}};
\node[op=deccolor, above=0.08cm of d1.north] {Lin + BN + LReLU + Drop};

\node[layer, fill=deccolor!15, right=1.5cm of d1, text width=1.6cm] (d2) {\textbf{256}};
\node[op=deccolor, above=0.08cm of d2.north] {Lin + BN + LReLU + Drop};

\node[layer, fill=deccolor!20, right=1.5cm of d2, text width=1.6cm] (d3) {\textbf{512}};
\node[op=deccolor, above=0.08cm of d3.north] {Lin + BN + LReLU + Drop};

\node[layer, fill=deccolor!8, right=1.8cm of d3, text width=1.6cm] (output) {\textbf{Output}\\$n_{genes}$};

% Arrows
\draw[arr, enccolor] (input) -- (e1);
\draw[arr, enccolor] (e1) -- (e2);
\draw[arr, enccolor] (e2) -- (e3);
\draw[arr, enccolor] (e3) -- (z);
\draw[arr, deccolor] (z) -- (d1);
\draw[arr, deccolor] (d1) -- (d2);
\draw[arr, deccolor] (d2) -- (d3);
\draw[arr, deccolor] (d3) -- (output);

% Braces
\draw[brace, enccolor] (input.south west) -- (e3.south east) node[midway, below=12pt, font=\small\bfseries, text=enccolor] {Encoder};
\draw[brace, deccolor] (d1.south west) -- (output.south east) node[midway, below=12pt, font=\small\bfseries, text=deccolor] {Decoder (symmetric)};

\end{tikzpicture}%
}
\caption{Architecture scRAW. Hidden layers : Linear $\to$ BatchNorm $\to$ LeakyReLU $\to$ Dropout(0.1). Latent space : pure linear projection (128D). Batch conditioning d\'esactiv\'e.}
\label{fig:architecture}
\end{figure}

\begin{parambox}{Param\`etres network}
\begin{tabular}{lll}
\texttt{z\_dim} & \textbf{128} & Dimension du latent space \\
\texttt{encoder\_layers} & \textbf{512, 256, 128} & Hidden layers de l'encoder \\
\texttt{decoder\_layers} & \textbf{(sym\'etrique)} & Decoder mirror : 128, 256, 512 \\
\texttt{activation} & \textbf{leaky\_relu} & Activation function \\
\texttt{dropout} & \textbf{0.1} & Dropout rate \\
\end{tabular}
\end{parambox}

\subsection{Adversarial discriminator (DANN) --- optionnel}
\label{sec:dann_archi}

Le module DANN (Domain-Adversarial Neural Network) permet la correction d'effets de batch en configuration multi-batch. Il ajoute un discriminateur adversarial sur l'espace latent qui force l'encoder \`a produire des repr\'esentations invariantes au batch d'origine.

\textbf{Architecture du discriminateur.} Le discriminateur est un MLP \`a deux couches appliqu\'e sur l'espace latent $z$ avec \textbf{gradient reversal} :
\begin{itemize}[nosep]
  \item Couche 1 : Linear($z_{dim}$, $z_{dim}/2$) + ReLU + Dropout(0.1)
  \item Couche 2 : Linear($z_{dim}/2$, $n_{batches}$)
  \item La perte est une cross-entropy sur la pr\'ediction du batch d'origine.
\end{itemize}
Le gradient reversal inverse le signe du gradient lors de la backpropagation ($-\lambda$), for\c{c}ant l'encoder \`a produire des repr\'esentations qui ne permettent pas de distinguer les batches.

\begin{parambox}{Param\`etres DANN}
\begin{tabular}{lll}
\texttt{use\_batch\_conditioning} & \textbf{false} & Conditionner le decoder sur le batch \\
\texttt{adversarial\_batch\_weight} & \textbf{0.0} & Poids de la perte adversariale (0 = d\'esactiv\'e) \\
\texttt{adversarial\_lambda} & \textbf{1.0} & \'Echelle du gradient reversal \\
\texttt{adversarial\_start\_epoch} & \textbf{0} & \'Epoque de d\'emarrage du DANN \\
\texttt{adversarial\_ramp\_epochs} & \textbf{0} & Dur\'ee du ramp-up lin\'eaire du DANN \\
\texttt{mmd\_batch\_weight} & \textbf{0.0} & Poids MMD (0 = d\'esactiv\'e) \\
\end{tabular}
\end{parambox}

% =============================================================================
\section{Entr\'ee des donn\'ees et reconstruction}
\label{sec:reconstruction}
% =============================================================================

\subsection{Double flux de donn\'ees}

scRAW maintient \textbf{deux repr\'esentations} des donn\'ees en parall\`ele :

\begin{figure}[H]
\centering
\resizebox{\linewidth}{!}{%
\begin{tikzpicture}[
  node distance=0.8cm and 1.5cm,
  box/.style={rectangle, draw, thick, rounded corners=5pt, align=center,
    minimum width=3cm, minimum height=1.0cm, font=\footnotesize},
  arr/.style={-{Stealth[length=2.5mm]}, thick},
]

% Raw counts path
\node[box, fill=gray!10] (raw) {\textbf{Raw counts}\\$X_{counts} \in \mathbb{N}^{n \times G}$\\(\texttt{layers['original\_X']})};

\node[box, fill=enccolor!10, right=2.5cm of raw, text width=4.8cm] (transform) {\textbf{Transformation NB d'entr\'ee}\\$x_{ij}^{enc} = T_{\text{NB}}(x_{ij})$\\(par d\'efaut $T_{\text{NB}}=\log(1+\cdot)$)};

\node[box, fill=enccolor!15, right=2.5cm of transform] (encoder) {\textbf{Encodeur}\\$z = f_{enc}(X_{enc})$};

% Decoder path
\node[box, fill=deccolor!10, below=1.5cm of encoder] (decoder) {\textbf{D\'ecodeur}\\$\text{logits} = f_{dec}(z)$};

% NB loss
\node[box, fill=nbblue!15, below=1.5cm of decoder, text width=4cm] (nbloss) {\textbf{Perte NB}\\$\mu = g(r_{ij}, s_i)$\\$\mathcal{L} = -\log P_{NB}(X_{counts} \mid \mu, \theta)$};

% Arrows
\draw[arr] (raw) -- node[above, font=\tiny]{\texttt{nb\_input\_transform}} (transform);
\draw[arr] (transform) -- (encoder);
\draw[arr] (encoder) -- (decoder);
\draw[arr] (decoder) -- (nbloss);
\draw[arr, dashed, gray] (raw) |- node[near end, above, font=\tiny]{Original counts (target)} (nbloss);

% Size factor
\node[box, fill=gray!8, left=1.5cm of nbloss, text width=2.5cm] (sf) {\textbf{Facteur de taille}\\$s_i = \mathrm{clip}\!\left(\frac{\text{lib\_size}_i}{\mathrm{med}_{+}},10^{-3},10^3\right)$};
\draw[arr, gray] (sf) -- (nbloss);

\end{tikzpicture}%
}
\caption{Double flux de donn\'ees (branche NB). Avec \texttt{original\_X} disponible, l'encoder re\c{c}oit $T_{\text{NB}}(X_{counts})$ (par d\'efaut \texttt{log1p} des raw counts), tandis que la perte NB est calcul\'ee contre les raw counts originaux.}
\label{fig:double_flux}
\end{figure}

\subsection{Transformation de l'entr\'ee encodeur en mode NB}

\textbf{Principe (impl\'ementation r\'eelle).} En mode de reconstruction NB, si \texttt{adata.layers['original\_X']} est disponible, l'algorithme reconstruit explicitement l'entr\'ee mod\`ele \`a partir de cette couche et ne r\'eutilise pas directement \texttt{adata.X} :
\[
  X_{enc} = T_{\text{NB}}(X_{counts}), \qquad T_{\text{NB}} \in \{\mathrm{identity},\log(1+\cdot)\}
\]
\textit{Description des termes :}
\begin{itemize}[nosep]
  \item $X_{counts}$ : matrice issue de \texttt{adata.layers['original\_X']}.
  \item $X_{enc}$ : matrice effectivement inject\'ee \`a l'encodeur.
  \item $T_{\text{NB}}$ : transformation contr\^ol\'ee par \texttt{nb\_input\_transform}.
\end{itemize}

\begin{enumerate}[nosep]
  \item \textbf{Cas par d\'efaut} (\texttt{nb\_input\_transform=log1p}) :
  \[
    x^{enc}_{ij} = \log(1 + x_{ij})
  \]
  \textit{Description des termes :}
  \begin{itemize}[nosep]
    \item $x_{ij}$ : count brut du g\`ene $j$ pour la cellule $i$.
    \item $x^{enc}_{ij}$ : valeur inject\'ee \`a l'encodeur en mode NB.
  \end{itemize}
  \item \textbf{Autres options} :
  \[
    x^{enc}_{ij}=x_{ij} \;\; (\texttt{none/identity})
  \]
\end{enumerate}

\textbf{Lien avec le reste du pipeline.}
\begin{itemize}[nosep]
  \item Le preprocessing global (\texttt{normalize\_total}, \texttt{log1p}, HVG, scale) produit \texttt{adata.X}, mais en branche NB l'entr\'ee encodeur est recalcul\'ee depuis \texttt{original\_X} comme ci-dessus.
  \item La perte NB est calcul\'ee contre les \textbf{raw counts} de \texttt{original\_X} quand cette couche existe ; sinon, fallback sur \texttt{adata.X}.
\end{itemize}

\subsection{Distribution Negative Binomial}

Les donn\'ees scRNA-seq pr\'esentent une \textbf{sur-dispersion} : la variance des counts est bien plus grande que leur moyenne, contrairement \`a une distribution de Poisson (o\`u $\text{Var}=\mu$). La distribution Negative Binomial mod\'elise cela explicitement via un param\`etre de dispersion $\theta$ :
\[
  \text{Var}_{NB} = \mu + \frac{\mu^2}{\theta}
\]
\textit{Description des termes :}
\begin{itemize}[nosep]
  \item $\mu$ : moyenne de comptage pr\'edite par le mod\`ele.
  \item $\theta$ : param\`etre de dispersion NB (\texttt{nb\_theta}, configurable).
\end{itemize}
Plus $\theta$ est petit, plus la sur-dispersion est importante ; $\theta \to \infty$ redonne une Poisson. Par d\'efaut, \texttt{nb\_theta=10.0}, mais cette valeur reste ajustable.

La perte de reconstruction utilise la \textbf{log-vraisemblance n\'egative NB} avec les m\^emes garde-fous num\'eriques que le code.

\textbf{Facteurs de taille (impl\'ementation exacte)} :
\[
  \mathrm{lib}_i = \sum_{j=1}^{G} x_{ij}, \qquad
  \mathrm{med}_{+} = \mathrm{m\'ediane}\!\left(\{\mathrm{lib}_u \mid \mathrm{lib}_u>0\}\right)
\]
\[
  s_i = \mathrm{clip}\!\left(\frac{\mathrm{lib}_i}{\max(\mathrm{med}_{+},10^{-6})},\,10^{-3},\,10^{3}\right)
\]
\textit{Description des termes :}
\begin{itemize}[nosep]
  \item $\mathrm{lib}_i$ : taille de biblioth\`eque brute de la cellule $i$.
  \item $\mathrm{med}_{+}$ : m\'ediane des tailles de biblioth\`eque strictement positives (fallback implicite \`a 1 si l'ensemble est vide).
  \item $s_i$ : size factor utilis\'e dans la param\'etrisation NB.
\end{itemize}

\textbf{Param\'etrisation de $\mu$} :
\[
  \theta'=\max(\theta,10^{-4}), \qquad x'_{ij}=\max(x_{ij},0)
\]
\[
  \mu_{ij}=
  \begin{cases}
    \mathrm{clip}\!\left(\exp\!\left(\mathrm{clip}(r_{ij},-12,12)\right)\cdot s_i,\;10^{-8},\;\mu_{\max}\right), & \text{si } s_i \text{ est disponible}\\
    \mathrm{softplus}(r_{ij}) + 10^{-4}, & \text{sinon}
  \end{cases}
\]
\textit{Description des termes :}
\begin{itemize}[nosep]
  \item $r_{ij}$ : sortie brute (\textit{logit}) du d\'ecodeur pour la cellule $i$, g\`ene $j$.
  \item $\mu_{\max}$ : borne sup\'erieure \texttt{nb\_mu\_clip\_max}.
  \item $\theta'$ et $x'_{ij}$ : versions clamp\'ees utilis\'ees dans le calcul num\'erique.
\end{itemize}

\textbf{NLL \'el\'ementaire (par cellule et par g\`ene)} :
pour chaque cellule $i$ et g\`ene $j$, le code calcule une perte NB locale $\ell^{NB}_{ij}$ :

\[
  \ell^{NB}_{ij} = -\Big[
  \ln\Gamma(x'_{ij}+\theta')-\ln\Gamma(\theta')-\ln\Gamma(x'_{ij}+1)
  + \theta' \big(\ln(\theta'+\varepsilon)-\ln(\theta'+\mu_{ij}+\varepsilon)\big)
  + x'_{ij}\big(\ln(\mu_{ij}+\varepsilon)-\ln(\theta'+\mu_{ij}+\varepsilon)\big)
  \Big]
\]
\[
  \ell^{NB}_{ij} \leftarrow \mathrm{nan\_to\_num}\!\left(\ell^{NB}_{ij},\;\mathrm{nan}=10^{3},\;+\infty=10^{3},\;-\infty=10^{3}\right),\quad \varepsilon=10^{-8}
\]

Sans masque, la perte de reconstruction par cellule est ensuite :
\[
  \mathcal{L}^{NB}_i = \frac{1}{G}\sum_{j=1}^{G}\ell^{NB}_{ij}
\]
Avec masque, l'agr\'egation suit la Section~\ref{sec:reconstruction} (paragraphe ``Agr\'egation de la perte avec masque'').

o\`u :
\begin{itemize}[nosep]
  \item $x_{ij}$ : raw count original du g\`ene $j$ dans la cellule $i$.
  \item $\theta$ : dispersion NB globale (hyperparam\`etre \texttt{nb\_theta}, non apprise).
  \item $\ln\Gamma(\cdot)$ : log-factorielle g\'en\'eralis\'ee.
\end{itemize}

\subsection{Reconstruction MSE (mode \texttt{reconstruction\_distribution=mse})}

\textbf{Ce que fait le code.} Si \texttt{reconstruction\_distribution} n'est pas \texttt{nb}, la branche MSE est utilis\'ee :
\[
  X_{enc}=X_{proc}=\texttt{adata.X}, \qquad X_{tgt}=X_{proc}
\]
\textit{Description des termes :}
\begin{itemize}[nosep]
  \item $X_{proc}$ : matrice pr\'etrait\'ee (sortie preprocessing Scanpy + scaling).
  \item $X_{enc}$ : entr\'ee de l'encodeur en mode MSE.
  \item $X_{tgt}$ : cible de reconstruction en mode MSE.
\end{itemize}

\textbf{Perte MSE \'el\'ementaire} (par cellule $i$, g\`ene $j$) :
\[
  \ell^{MSE}_{ij} = \left(\hat{x}_{ij} - x_{ij}\right)^2
\]
Sans masque :
\[
  \mathcal{L}^{MSE}_i = \frac{1}{G}\sum_{j=1}^{G}\ell^{MSE}_{ij}
\]
Avec masque, l'agr\'egation est identique \`a la NB (Section~\ref{sec:reconstruction}, paragraphe ``Agr\'egation de la perte avec masque'').

\textbf{Lien avec les phases d'entra\^inement} :
\[
  \mathcal{L}_{warmup}=\frac{1}{N}\sum_{i=1}^{N}\mathcal{L}^{MSE}_i
  \qquad\text{(poids uniformes)}
\]
\[
  \mathcal{L}_{recon,weighted}=\frac{1}{N}\sum_{i=1}^{N}w_i\,\mathcal{L}^{MSE}_i
  \qquad\text{(phase pond\'er\'ee)}
\]
Ce comportement est strictement le m\^eme sch\'ema que la NB, en rempla\c{c}ant la NLL NB par la MSE.

\subsection{Masquage de g\`enes (style scMAE)}

Le masquage est un \textbf{bruit d'entr\'ee stochastique} appliqu\'e pendant le \textbf{warm-up uniquement}. Concr\`etement, \`a \textbf{chaque mini-batch} (donc plusieurs fois par \'epoque), un nouveau masque est tir\'e ind\'ependamment pour chaque couple cellule--g\`ene ; il n'y a pas de masque fixe conserv\'e entre les \'epoques.

\[
M_{ij}^{(e,b)} \sim \mathrm{Bernoulli}(p), \quad p=\texttt{masking\_rate}
\]
\[
\tilde{x}_{ij}^{(e,b)} =
\begin{cases}
v_{\text{mask}}, & \text{si } M_{ij}^{(e,b)}=1 \\
x_{ij}, & \text{sinon}
\end{cases}
\]
\textit{Description des termes :}
\begin{itemize}[nosep]
  \item $M_{ij}^{(e,b)}$ : variable binaire de masquage pour la cellule $i$, g\`ene $j$, \`a l'\'epoque $e$ et mini-batch $b$.
  \item $p$ : probabilit\'e de masquage (\texttt{masking\_rate}).
  \item $v_{\text{mask}}$ : valeur inject\'ee aux positions masqu\'ees (\texttt{masking\_value}).
  \item $\tilde{x}_{ij}^{(e,b)}$ : entr\'ee corrompue transmise \`a l'encoder.
\end{itemize}

L'autoencoder re\c{c}oit donc $\tilde{x}$ en entr\'ee et doit reconstruire la cellule originale. Cela agit comme un d\'ebruitage (type scMAE) et limite la m\'emorisation.

\paragraph{Agr\'egation de la perte avec masque}
Quand un masque est pr\'esent, la perte par g\`ene $\ell_{ij}$ (NB ou MSE) est agr\'eg\'ee avec :
\[
a_{ij} = M_{ij}\cdot\omega + (1-M_{ij})\cdot(1-\omega), \quad \omega=\texttt{masked\_recon\_weight}
\]
\[
\mathcal{L}_i = \frac{\sum_j a_{ij}\,\ell_{ij}}{\sum_j a_{ij}}
\]
\textit{Description des termes :}
\begin{itemize}[nosep]
  \item $a_{ij}$ : coefficient de pond\'eration appliqu\'e \`a la perte du g\`ene $j$ pour la cellule $i$.
  \item $\omega$ : poids relatif des positions masqu\'ees (\texttt{masked\_recon\_weight}).
  \item $\ell_{ij}$ : perte \'el\'ementaire par g\`ene (NB ou MSE).
  \item $\mathcal{L}_i$ : perte de reconstruction agr\'eg\'ee pour la cellule $i$.
\end{itemize}
Avec $\omega>0.5$, les positions masqu\'ees contribuent davantage que les positions non masqu\'ees.

\paragraph{Chronologie du masquage}
\begin{itemize}[nosep]
  \item Warm-up (\'epoques 0--29) : masquage actif, r\'e\'echantillonn\'e \`a chaque mini-batch.
  \item Phase pond\'er\'ee (\'epoques 30--119) : masquage \textbf{d\'esactiv\'e} car \texttt{masking\_apply\_weighted=false}.
\end{itemize}

\begin{figure}[H]
\centering
\begin{tikzpicture}[
  gene/.style={rectangle, minimum width=0.35cm, minimum height=0.7cm, draw=gray!50, thick, font=\tiny},
]

% Original
\node[font=\footnotesize\bfseries] at (-1.2, 0) {Original :};
\foreach \i/\v/\c in {0/3.2/enccolor!20, 1/0.0/gray!10, 2/1.7/enccolor!20, 3/0.0/gray!10, 4/5.1/enccolor!30, 5/0.8/enccolor!10, 6/0.0/gray!10, 7/2.3/enccolor!20, 8/4.6/enccolor!25, 9/0.1/enccolor!5, 10/0.0/gray!10, 11/3.8/enccolor!25} {
  \node[gene, fill=\c] at (\i*0.45, 0) {\v};
}

% Masked
\node[font=\footnotesize\bfseries] at (-1.2, -1.2) {Masqu\'e :};
\foreach \i/\v/\c/\m in {0/3.2/enccolor!20/0, 1/0.0/gray!10/0, 2/\textcolor{dannred}{0}/dannred!15/1, 3/0.0/gray!10/0, 4/5.1/enccolor!30/0, 5/\textcolor{dannred}{0}/dannred!15/1, 6/0.0/gray!10/0, 7/2.3/enccolor!20/0, 8/\textcolor{dannred}{0}/dannred!15/1, 9/0.1/enccolor!5/0, 10/0.0/gray!10/0, 11/3.8/enccolor!25/0} {
  \node[gene, fill=\c] at (\i*0.45, -1.2) {\v};
}

% Brace for masked
\node[font=\tiny, text=dannred] at (6.2, -1.2) {$\approx 20\%$};

\end{tikzpicture}
\caption{Masquage al\'eatoire de 20\% des g\`enes en entr\'ee. Les positions masqu\'ees (en rouge) sont remplac\'ees par 0 et pond\'er\'ees avec \texttt{masked\_recon\_weight=0.75} (positions non masqu\'ees : 0.25) dans la perte de reconstruction.}
\label{fig:masking}
\end{figure}

\begin{parambox}{Param\`etres reconstruction --- valeurs par d\'efaut}
\begin{tabular}{lll}
\texttt{reconstruction\_distribution} & \textbf{nb} & Distribution Negative Binomial \\
\texttt{nb\_theta} & \textbf{10.0} & Dispersion globale (fixe, non apprise) \\
\texttt{nb\_input\_transform} & \textbf{log1p} & Transformation de l'entr\'ee encoder \\
\texttt{masking\_rate} & \textbf{0.2} & 20\% des g\`enes masqu\'es (warm-up seul) \\
\texttt{masked\_recon\_weight} & \textbf{0.75} & Poids relatif des positions masqu\'ees dans la perte \\
\texttt{masking\_apply\_weighted} & \textbf{false} & Masquage d\'esactiv\'e en phase 2 \\
\end{tabular}
\end{parambox}

% =============================================================================
\section{D\'eroulement de l'entra\^inement}
\label{sec:phases}
% =============================================================================

L'entra\^inement est organis\'e en deux phases distinctes afin d'initialiser les poids puis de raffiner la repr\'esentation des populations rares.

\begin{figure}[H]
\centering
\resizebox{\linewidth}{!}{%
\begin{tikzpicture}[
  x=0.14cm, y=1cm,
]

% Timeline axis (scaled: 1 unit = 1 epoch, displayed up to 120)
\draw[-{Stealth}, thick] (0,0) -- (130,0) node[right, font=\small] {\'Epoques};
\foreach \x in {0,10,20,30,40,50,60,70,80,90,100,110,120} {
  \draw[thick] (\x, -0.1) -- (\x, 0.1);
  \node[below, font=\tiny] at (\x, -0.15) {\x};
}

% Phase 1: Warm-up (0-29)
\fill[phaseblue!15] (0, 0.3) rectangle (30, 2.8);
\node[font=\small\bfseries, text=phaseblue, align=center] at (15, 2.3) {Phase 1 : Warm-up};
\node[font=\tiny, align=center, text width=3.5cm] at (15, 1.5) {Perte NB uniforme\\Masquage 20\% actif\\Poids $w_i = 1$ pour tous};

% Phase 2: Weighted (30-119)
\fill[phaseorange!12] (30, 0.3) rectangle (120, 2.8);
\node[font=\small\bfseries, text=phaseorange, align=center] at (75, 2.3) {Phase 2 : Pond\'er\'ee + R\'egularisation rare};
\node[font=\tiny, align=center, text width=10cm] at (75, 1.5) {Perte NB pond\'er\'ee : $\mathcal{L} = \sum w_i \cdot \mathcal{L}_{NB}$\\Masquage d\'esactiv\'e \quad Poids dynamiques $w_i \in [0.25, 10]$};

% Triplet ramp (start epoch 35, ramp 20 ep)
\fill[tripletyellow!15] (35, -0.4) rectangle (120, -1.2);
\draw[tripletyellow!80!black, thick, -{Stealth}] (35, -0.8) -- (55, -0.8) node[midway, above, font=\tiny] {Ramp-up $T_{\text{ramp}}$ \'ep.};
\draw[tripletyellow!80!black, thick] (55, -0.8) -- (120, -0.8);
\node[font=\tiny, text=tripletyellow!80!black] at (80, -1.05) {Triplet Loss active ($\lambda_{rare}$, d\`es \'ep.~$e_{\text{start}}$)};

% Weight updates markers
\foreach \e in {40, 50, 60, 70, 80, 90, 100, 110} {
  \draw[weightpurple, thick] (\e, 0.3) -- (\e, 0.7);
  \node[font=\tiny, text=weightpurple, rotate=90, anchor=west] at (\e, 0.7) {$w$};
}

% Transition marker (epoch 30)
\draw[thick, dashed, gray] (30, -1.5) -- (30, 3.0);
\node[font=\tiny\bfseries, gray, rotate=90, anchor=south] at (30, 3.2) {Transition};

% Triplet start marker (epoch 35)
\draw[thick, dashed, tripletyellow!80!black] (35, -1.5) -- (35, 3.0);
\node[font=\tiny\bfseries, text=tripletyellow!80!black, rotate=90, anchor=south] at (35, 3.2) {Triplet start};

% Legend
\node[font=\tiny, weightpurple] at (112, 0.95) {$\uparrow$ MAJ poids};

\end{tikzpicture}%
}
\caption{Chronologie de l'entra\^inement scRAW (exemple : 120 \'epoques, warm-up 30 + phase pond\'er\'ee 90). Les poids dynamiques sont initialis\'es \`a la transition puis mis \`a jour p\'eriodiquement. La Triplet Loss s'active apr\`es le warm-up avec un ramp-up lin\'eaire.}
\label{fig:timeline}
\end{figure}

% -----------------------------------------------------------------------------
\subsection{Phase 1 : Warm-up (\'epoques 0 \`a 29)}
\label{sec:phase1}

\textbf{Objectif} : initialiser l'encoder/d\'ecodeur avec une repr\'esentation stable, sans biais de pond\'eration.

\textbf{Fonction de perte} :
\[
  \mathcal{L}_{warmup} = \frac{1}{N} \sum_{i=1}^{N} \mathcal{L}_{recon}(x_i, \hat{x}_i)
\]
\textit{Description des termes :}
\begin{itemize}[nosep]
  \item $N$ : nombre de cellules dans le batch global consid\'er\'e pour l'\'epoque.
  \item $x_i$ : profil d'entr\'ee de la cellule $i$.
  \item $\hat{x}_i$ : reconstruction produite par le d\'ecodeur.
  \item $\mathcal{L}_{recon}$ : perte de reconstruction par cellule (NB ou MSE selon configuration).
\end{itemize}

Toutes les cellules sont trait\'ees de mani\`ere \'egale (poids uniforme $w_i = 1$). Le masquage est r\'e\'echantillonn\'e \`a chaque mini-batch et force l'autoencoder \`a g\'en\'eraliser plut\^ot que m\'emoriser. L'optimiseur est Adam avec un scheduler \textbf{CosineAnnealingLR} ($\eta_{min} = \texttt{lr} \times 0.01$) et un \textbf{gradient clipping} par norme (\texttt{max\_norm}=5.0).

\begin{parambox}{Param\`etres warm-up}
\begin{tabular}{lll}
\texttt{warmup\_epochs} & \textbf{30} & Dur\'ee de la phase warm-up \\
\texttt{batch\_size} & \textbf{256} & Taille des mini-batches \\
\texttt{lr} & \textbf{0.001} & Taux d'apprentissage (Adam) \\
\texttt{scheduler} & \textbf{CosineAnnealingLR} & Avec $\eta_{min} = \texttt{lr} \times 0.01$ \\
\texttt{gradient clipping} & \textbf{max\_norm=5.0} & Clip des gradients par norme \\
\end{tabular}
\end{parambox}

% -----------------------------------------------------------------------------
\subsection{Transition : calcul initial des poids}
\label{sec:transition}

\`A la fin du warm-up (\'epoque 30), l'entra\^inement s'interrompt bri\`evement pour :
\begin{enumerate}[nosep]
  \item \textbf{Encoder} tout le jeu de donn\'ees dans l'espace latent courant.
  \item \textbf{Fixer un $K$ de travail pour le pseudo-labeling} :
  \begin{itemize}[nosep]
    \item mode oracle : $K$ fourni (nombre de types cellulaires connus, via \texttt{n\_clusters\_effective}),
    \item vrai non supervis\'e : $K$ auto-estim\'e depuis les embeddings warm-up.
  \end{itemize}
  \item \textbf{Pseudo-labeling} (ici Leiden) sur l'espace latent.
  \item \textbf{Calculer les poids} bicrit\`eres (fr\'equence + densit\'e) pour chaque cellule.
\end{enumerate}

Ce processus est r\'ep\'et\'e tous les 10 cycles pendant la phase 2 (voir Section~\ref{sec:weighting}).

% -----------------------------------------------------------------------------
\subsection{Phase 2 : Entra\^inement pond\'er\'e (\'epoques 30 \`a 120)}
\label{sec:phase2}

\textbf{Fonction de perte totale} :
\[
  \boxed{\mathcal{L}_{total} = \underbrace{\frac{1}{N}\sum_{i=1}^{N} w_i \cdot \mathcal{L}_{recon}(x_i, \hat{x}_i)}_{\text{Reconstruction pond\'er\'ee}} \;+\; \underbrace{\rho(e) \cdot \lambda_{rare} \cdot \mathcal{L}_{triplet}}_{\text{R\'egularisation rare}}}
\]
\textit{Description des termes :}
\begin{itemize}[nosep]
  \item $\mathcal{L}_{total}$ : objectif optimis\'e en phase pond\'er\'ee.
  \item $N$ : nombre de cellules.
  \item $w_i$ : poids de la cellule $i$ issu du syst\`eme de pond\'eration (Section~\ref{sec:weighting}).
  \item $\mathcal{L}_{recon}(x_i,\hat{x}_i)$ : perte de reconstruction par cellule (NB si \texttt{reconstruction\_distribution=nb}, MSE si \texttt{reconstruction\_distribution=mse}).
  \item $\rho(e)$ : coefficient de ramp-up \`a l'\'epoque $e$.
  \item $\lambda_{rare}$ : poids de la Triplet Loss.
  \item $\mathcal{L}_{triplet}$ : terme de r\'egularisation latente pour les cellules rares.
\end{itemize}

o\`u :
\begin{itemize}[nosep]
  \item $w_i \in [w_{\min}, w_{\max}]$ est le poids dynamique bicrit\`ere de la cellule $i$ (\texttt{min\_cell\_weight}, \texttt{max\_cell\_weight}).
  \item $\rho(e) = \max\!\left(0,\; \min\!\left(1, \dfrac{e - e_{\text{start}}}{T_{\text{ramp}}}\right)\right)$ est le coefficient de ramp-up lin\'eaire, avec $e_{\text{start}}=\texttt{rare\_triplet\_start\_epoch}$ et $T_{\text{ramp}}$ la dur\'ee de mont\'ee.
  \item $\lambda_{rare}$ est le poids de la r\'egularisation rare (\texttt{rare\_triplet\_weight}).
\end{itemize}

\textbf{Note} : en configuration mono-batch, la correction DANN est d\'esactiv\'ee (\texttt{adversarial\_batch\_weight=0.0}) et la perte totale se r\'eduit \`a reconstruction pond\'er\'ee + Triplet Loss. En multi-batch, un terme adversarial $\lambda_{adv} \cdot \mathcal{L}_{DANN}$ s'ajoute.

\textbf{Changement cl\'e} : le masquage est \textbf{d\'esactiv\'e} en phase 2 (\texttt{masking\_apply\_weighted = false}) pour ne pas aveugler les signaux d\'elicats des cellules rares, dont chaque g\`ene compte.

\begin{parambox}{Param\`etres entra\^inement global --- valeurs par d\'efaut}
\begin{tabular}{lll}
\texttt{epochs} & \textbf{120} & Nombre total d'\'epoques (warm-up + pond\'er\'ee) \\
\texttt{warmup\_epochs} & \textbf{30} & Dur\'ee de la phase 1 (\'epoques 0--29) \\
\texttt{batch\_size} & \textbf{256} & Cellules par mini-batch \\
\texttt{lr} & \textbf{0.001} & Taux d'apprentissage Adam \\
\texttt{scheduler} & \textbf{CosineAnnealingLR} & $\eta_{min} = \texttt{lr} \times 0.01$ \\
\texttt{gradient clipping} & \textbf{max\_norm=5.0} & Clip des gradients par norme \\
\texttt{random\_state} & \textbf{42} & Graine pour reproductibilit\'e \\
\end{tabular}
\end{parambox}

% =============================================================================
\section{Pond\'eration dynamique bicrit\`ere}
\label{sec:weighting}
% =============================================================================

Le syst\`eme de pond\'eration attribue \`a chaque cellule un poids $w_i$ calcul\'e en combinant deux mesures de raret\'e.

\begin{figure}[H]
\centering
\resizebox{0.75\linewidth}{!}{%
\begin{tikzpicture}[
  node distance=0.8cm and 1.0cm,
  comp/.style={rectangle, draw, thick, rounded corners=6pt, align=center,
    text width=4.3cm, minimum height=1.6cm, font=\scriptsize},
  op/.style={circle, draw, thick, fill=white, minimum size=0.8cm, font=\small\bfseries},
  arr/.style={-{Stealth[length=2.5mm]}, thick},
]

% Input: latent space
\node[comp, fill=latentcolor!15, text width=3.5cm] (latent) {\textbf{Espace latent $z$}\\(toutes les cellules)};

% Two branches
\node[comp, fill=weightpurple!10, below left=1.3cm and -0.1cm of latent] (freq) {%
  \textbf{Poids de fr\'equence}\\[4pt]
  Leiden $\to$ pseudo-labels\\[6pt]
  $w_{freq} = \left(\frac{1}{\text{freq}_c}\right)^{\alpha}$\\[6pt]
  Normalis\'e : $\overline{w} = 1$};

\node[comp, fill=weightpurple!10, below right=1.3cm and -0.1cm of latent] (density) {%
  \textbf{Poids de densit\'e}\\[4pt]
  Distance $k$-NN ($k=\texttt{density\_knn\_k}$)\\[8pt]
  $w_{dens} = \left(\frac{d_i^{(k)}}{\text{m\'ediane}}\right)^{\beta}$\\[6pt]
  Clip : $[0.05, 5.0]$};

% Fusion
\node[op, below=2.0cm of latent, fill=weightpurple!20] (fusion) {$+$};

% Result
\node[comp, fill=weightpurple!15, below=1.0cm of fusion, text width=5cm] (result) {%
  \textbf{Poids final}\\[2pt]
  $w_i = \gamma \cdot w_{freq} + (1-\gamma) \cdot w_{dens}$\\[2pt]
  Momentum : $w^{(t)} = \mu \cdot w^{(t-1)} + (1-\mu) \cdot w^{fresh}$\\[2pt]
  Clip final : $[0.25,\; 10.0]$};

% Arrows
\draw[arr] (latent) -| (freq);
\draw[arr] (latent) -| (density);
\draw[arr] (freq) |- (fusion) node[near end, above left, font=\tiny\bfseries]{$\gamma$};
\draw[arr] (density) |- (fusion) node[near end, above right, font=\tiny\bfseries]{$1-\gamma$};
\draw[arr, weightpurple] (fusion) -- (result);

% Application arrow
\node[comp, fill=phaseorange!10, below=0.8cm of result, text width=5cm] (apply) {%
  \textbf{Application dans la perte}\\[2pt]
  $\mathcal{L}_{recon} = \sum w_i \cdot \mathcal{L}_{NB}(x_i, \hat{x}_i)$};
\draw[arr, phaseorange] (result) -- (apply);

\end{tikzpicture}%
}
\caption{Syst\`eme de pond\'eration bicrit\`ere. Les deux composantes (fr\'equence de cluster et densit\'e locale) sont combin\'ees additivement via $\gamma$, puis stabilis\'ees par momentum EMA $\mu$ et clipp\'ees \`a $[0.25, 10]$.}
\label{fig:weighting}
\end{figure}

\subsection{Composante 1 : fr\'equence de pseudo-cluster}

\textbf{Principe.} Leiden est appliqu\'e sur l'espace latent courant pour produire des pseudo-labels $\{c_1, \ldots, c_N\}$ (pas les vrais labels biologiques). Ces pseudo-clusters d\'efinissent des groupes temporaires utilis\'es uniquement pour le calcul des poids --- ils n'interviennent pas dans l'\'evaluation finale.

\begin{enumerate}[nosep]
  \item \textbf{Leiden} sur les $k$-NN du latent space $z$ $\to$ pseudo-labels $c_i \in \{1,\ldots,K\}$.
  \item Fr\'equence relative de chaque cluster : $\text{freq}_c = |c| / N$ (proportion de cellules dans $c$).
  \item Poids inverse-fr\'equence soft :
  \[
    w_{freq,i} = \left(\frac{1}{\mathrm{freq}_{c_i}}\right)^{\alpha}
  \]
  \textit{Description des termes :}
  \begin{itemize}[nosep]
    \item $w_{freq,i}$ : composante fr\'equence du poids de la cellule $i$.
    \item $\mathrm{freq}_{c_i}$ : fr\'equence relative du pseudo-cluster auquel appartient la cellule $i$.
    \item $\alpha$ : exposant de la composante fr\'equence (\texttt{weight\_exponent}).
  \end{itemize}
  L'exposant $\alpha < 1$ mod\`ere l'amplification : sans lui ($\alpha=1$), un cluster \`a 0.1\% des cellules recevrait un poids $\times 1000$ --- instable.
  \item Normalisation globale : $w_{freq,i} \leftarrow w_{freq,i} / \overline{w_{freq}}$, de sorte que \textbf{la moyenne des poids de fr\'equence sur tout le dataset soit 1}. Cela pr\'eserve l'\'echelle globale de la perte de reconstruction.
\end{enumerate}

\textbf{Exemple illustratif} (jeu de donn\'ees avec des populations de tailles vari\'ees) :

\begin{center}
\resizebox{\linewidth}{!}{%
\begin{tikzpicture}[x=0.08cm, y=1cm]
% Bars
\foreach \name/\n/\w/\xpos in {%
  beta/2525/0.57/0,
  alpha/2326/0.58/18,
  ductal/1077/0.73/36,
  acinar/958/0.76/54,
  delta/601/0.87/72,
  stellate/284/1.09/90,
  gamma/255/1.13/108,
  endoth./252/1.14/126,
  quiesc./173/1.27/144,
  macro./55/1.72/162,
  mast/25/2.18/180,
  epsilon/18/2.35/198,
  schwann/13/2.52/216,
  T~cell/7/2.85/234%
} {
  % Map weight [0.57, 2.85] → color intensity [20, 85]
  \pgfmathtruncatemacro{\clr}{min(85, max(20, int(\w * 30)))}
  \fill[weightpurple!\clr] (\xpos, 0) rectangle (\xpos+14, \w*0.35);
  \node[font=\tiny, rotate=45, anchor=west] at (\xpos+7, \w*0.35+0.05) {\name};
  \node[font=\tiny, weightpurple] at (\xpos+7, \w*0.35-0.08) {\w$\times$};
}

% Axis
\draw[-{Stealth}, thick] (-2, 0) -- (252, 0) node[right, font=\tiny] {Types cellulaires};
\draw[-{Stealth}, thick] (-2, 0) -- (-2, 1.2) node[above, font=\tiny] {$w_{freq}$};

\end{tikzpicture}%
}
\end{center}

\subsection{Composante 2 : densit\'e locale (kNN)}

\textbf{Principe.} Pour chaque cellule $i$, on mesure son degr\'e d'isolement dans l'espace latent via la distance \`a son $k^e$ plus proche voisin ($k=\texttt{density\_knn\_k}, born\'e \`a $[2,n-1]$) :
\[
  \tilde{w}^{dens}_i = \mathrm{clip}\!\left(\left(\frac{d_i^{(k)}}{\mathrm{m\'ediane}_{j}\bigl(d_j^{(k)}\bigr)}\right)^{\beta},\;0.05,\;c_{dens}\right),\qquad
  w_{dens,i}=\frac{\tilde{w}^{dens}_i}{\mathrm{moyenne}_{u}(\tilde{w}^{dens}_u)}
\]
\textit{Description des termes :}
\begin{itemize}[nosep]
  \item $w_{dens,i}$ : composante densit\'e du poids de la cellule $i$.
  \item $\tilde{w}^{dens}_i$ : poids densit\'e avant renormalisation de moyenne.
  \item $d_i^{(k)}$ : distance euclidienne entre la cellule $i$ et sa $k^e$ voisine la plus proche dans l'espace latent.
  \item $\mathrm{m\'ediane}_{j}(d_j^{(k)})$ : distance $k$-NN m\'ediane sur l'ensemble des cellules.
  \item $\beta$ : exposant de la composante densit\'e (\texttt{density\_weight\_exponent}).
  \item $c_{dens}$ : borne sup\'erieure de clipping (\texttt{density\_weight\_clip}, avec minimum effectif 0.05).
  \item $k$ : ordre de voisinage pour l'estimation de densit\'e (\texttt{density\_knn\_k}).
\end{itemize}

La normalisation par la m\'ediane rend le poids relatif \`a la densit\'e typique du dataset.

\begin{itemize}[nosep]
  \item \textbf{Cellule isol\'ee} (loin de ses voisins, population rare) : $d_i^{(k)} \gg \text{m\'ediane}$ $\Rightarrow$ $w_{dens,i} > 1$.
  \item \textbf{Cellule dense} (au centre d'un gros cluster) : $d_i^{(k)} \ll \text{m\'ediane}$ $\Rightarrow$ $w_{dens,i} < 1$.
  \item Clip pr\'e-fusion : $\tilde{w}^{dens}_i \in [0.05, c_{dens}]$ avant renormalisation, pour \'eviter les poids extr\^emes.
\end{itemize}

\textbf{Compl\'ementarit\'e avec la fr\'equence.} Une cellule peut appartenir \`a un petit pseudo-cluster Leiden (fort $w_{freq}$) tout en \'etant physiquement proche d'un gros cluster voisin (faible $w_{dens}$). La fusion additive via $\gamma$ \'equilibre ces deux informations et \'evite de sur-pond\'erer des cellules qui ne sont rares que sur le plan num\'erique mais pas morphologiquement isol\'ees.

\subsection{Mise \`a jour dynamique et stabilisation EMA}

P\'eriodiquement pendant la phase 2 (tous les \texttt{dynamic\_weight\_update\_interval} cycles), les poids sont enti\`erement recalcul\'es depuis l'espace latent courant : nouveau Leiden, nouvelles fr\'equences de clusters, nouvelle densit\'e $k$-NN. Ce recalcul n\'ecessite l'encodage de toutes les cellules mais est n\'ecessaire pour que les poids refl\`etent l'\'evolution de la repr\'esentation latente.

La fusion fr\'equence/densit\'e utilis\'ee juste avant l'EMA suit l'impl\'ementation :
\[
  c_i=\max(w_{freq,i},10^{-8})^{p_c},\qquad d_i=\max(w_{dens,i},10^{-8})^{p_d}
\]
\[
  w_i^{raw}=
  \begin{cases}
    c_i\cdot d_i, & \text{si }\texttt{weight\_fusion\_mode}=\texttt{multiplicative}\\
    \gamma c_i + (1-\gamma)d_i, & \text{si }\texttt{weight\_fusion\_mode}=\texttt{additive}
  \end{cases}
\]
\[
  w_i^{new}=\mathrm{clip}\!\left(\frac{w_i^{raw}}{\mathrm{moyenne}_{u}(w_u^{raw})},\,w_{\min},\,w_{\max}\right)
\]
\textit{Description des termes :}
\begin{itemize}[nosep]
  \item $p_c$, $p_d$ : puissances \texttt{cluster\_weight\_power} et \texttt{density\_weight\_power}.
  \item $\gamma$ : ratio de fusion \texttt{cluster\_density\_alpha} (appliqu\'e en mode additif).
  \item $w_{\min},w_{\max}$ : bornes de clipping final (\texttt{min\_cell\_weight}, \texttt{max\_cell\_weight}).
\end{itemize}

Pour \'eviter les \textbf{oscillations de gradient} li\'ees aux changements brutaux de poids, une EMA est appliqu\'ee pour les mises \`a jour suivantes ($t\ge 1$) :
\[
  \tilde{w}_i^{(t)} = \mu \cdot w_i^{(t-1)} + (1-\mu) \cdot w_i^{new}
\]
\[
  w_i^{(t)} = \mathrm{clip}\!\left(\frac{\tilde{w}_i^{(t)}}{\mathrm{moyenne}_{u}(\tilde{w}_u^{(t)})},\,w_{\min},\,w_{\max}\right)
\]
\textit{Description des termes :}
\begin{itemize}[nosep]
  \item $w_i^{(t)}$ : poids appliqu\'e \`a la cellule $i$ apr\`es lissage/renormalisation.
  \item $w_i^{new}$ : poids fra\^ichement recalcul\'e avant EMA.
  \item $w_i^{(t-1)}$ : poids utilis\'e au pas de mise \`a jour pr\'ec\'edent.
  \item $\mu$ : coefficient de momentum EMA (\texttt{dynamic\_weight\_momentum}).
\end{itemize}
Au premier calcul de phase pond\'er\'ee (transition warm-up $\to$ weighted), le code applique directement $w_i^{new}$ sans EMA.

\textbf{Chronologie des mises \`a jour :} initialisation \`a la transition (fin du warm-up), puis mises \`a jour r\'eguli\`eres. La premi\`ere mise \`a jour intervient apr\`es le d\'emarrage du weighting, une fois que la Triplet Loss commence \`a influencer l'espace latent.

\begin{parambox}{Param\`etres pond\'eration --- valeurs par d\'efaut}
\begin{tabular}{lll}
\texttt{pseudo\_label\_method} & \textbf{leiden} & M\'ethode de pseudo-labeling \\
\texttt{weight\_exponent} ($\alpha$) & \textbf{0.4} & Exposant inverse-fr\'equence \\
\texttt{density\_knn\_k} & \textbf{15} & Voisins pour densit\'e locale \\
\texttt{density\_weight\_exponent} ($\beta$) & \textbf{1.0} & Exposant distance kNN \\
\texttt{density\_weight\_clip} & \textbf{5.0} & Clip max densit\'e avant fusion \\
\texttt{weight\_fusion\_mode} & \textbf{additive} & $w = \gamma \cdot w_{freq} + (1-\gamma) \cdot w_{dens}$ \\
\texttt{cluster\_density\_alpha} ($\gamma$) & \textbf{0.6} & Ratio fr\'equence/densit\'e \\
\texttt{min\_cell\_weight} & \textbf{0.25} & Poids minimum apr\`es normalisation \\
\texttt{max\_cell\_weight} & \textbf{10.0} & Poids maximum \\
\texttt{dynamic\_weight\_update\_interval} & \textbf{10} & Recalcul tous les 10 cycles \\
\texttt{dynamic\_weight\_momentum} ($\mu$) & \textbf{0.7} & Inertie EMA \\
\end{tabular}
\end{parambox}

% =============================================================================
\section{R\'egularisation rare : Triplet Loss}
\label{sec:triplet}
% =============================================================================

La Triplet Loss structure \textbf{localement} l'espace latent autour des populations rares. Pour chaque cellule rare (ancre), elle attire un voisin du m\^eme pseudo-cluster (positif) tout en repoussant un intrus d'un autre cluster (n\'egatif), jusqu'\`a ce qu'une marge de s\'ecurit\'e soit respect\'ee.

\begin{figure}[H]
\centering
\resizebox{\linewidth}{!}{%
\begin{tikzpicture}[
  scale=1.2,
  cell/.style={circle, draw, thick, minimum size=0.5cm, inner sep=0, font=\tiny\bfseries},
]

% Background clusters
\fill[enccolor!8, rounded corners=10pt] (-0.5, -0.5) rectangle (2.5, 2.5);
\node[font=\tiny, enccolor, anchor=north west] at (-0.4, 2.4) {Cluster A (rare)};

\fill[deccolor!8, rounded corners=10pt] (4.5, -0.5) rectangle (7.5, 2.5);
\node[font=\tiny, deccolor, anchor=north west] at (4.6, 2.4) {Cluster B};

% Anchor
\node[cell, fill=tripletyellow!40, line width=2pt] (a) at (1.0, 1.0) {A};
\node[font=\tiny, below=0.35cm of a] {Ancre ($w \geq 1.2$)};

% Positive (hardest = farthest in same cluster)
\node[cell, fill=enccolor!30] (p) at (2.0, 2.0) {P};
\node[font=\tiny, right=0.15cm of p, text=enccolor] {Hard positive};

% Negative (semi-hard)
\node[cell, fill=deccolor!30] (n) at (4.8, 1.2) {N};
\node[font=\tiny, right=0.15cm of n, text=deccolor] {Semi-hard negative};

% Other cells
\foreach \x/\y in {0.3/0.5, 0.8/1.8, 1.5/0.3, 1.8/1.2} {
  \node[cell, fill=enccolor!15] at (\x, \y) {};
}
\foreach \x/\y in {5.2/0.5, 5.8/1.8, 6.5/0.8, 5.5/1.5, 6.2/1.2, 6.8/2.0} {
  \node[cell, fill=deccolor!15] at (\x, \y) {};
}

% Distances
\draw[enccolor, thick, <->] (a.north east) -- (p.south west) node[midway, sloped, above=3pt, font=\scriptsize] {$d_{pos}$ (max)};
\draw[deccolor, thick, <->] (a.east) -- (n.west) node[midway, above=3pt, font=\scriptsize] {$d_{neg}$ (semi-hard)};

% Margin annotation
\draw[thick, dashed, gray] (3.2, -0.8) -- (3.2, 3.2);
\node[font=\scriptsize, gray, align=center] at (3.2, 3.6) {Marge $m$\\entre clusters};

% Loss formula
\node[draw, thick, rounded corners=4pt, fill=tripletyellow!10, font=\footnotesize, align=center,
  text width=6cm] at (3.5, -1.5) {%
  $\mathcal{L}_{triplet} = \frac{1}{|\mathcal{A}|}\sum_{a} \max\big(0,\; d_{pos}(a) - d_{neg}(a) + m\big)$};

\end{tikzpicture}%
}
\caption{Triplet Loss avec semi-hard mining. L'ancre (cellule rare, $w \geq 1.2$) est attir\'ee vers son hard positive (alli\'e le plus \'eloign\'e) et repouss\'ee de son semi-hard negative (intrus le plus proche mais d\'ej\`a au-del\`a de $d_{pos}$). La marge $m$ garantit une s\'eparation minimale.}
\label{fig:triplet}
\end{figure}

\textbf{Strat\'egie de semi-hard mining} :
\begin{itemize}[nosep]
  \item \textbf{Hard positive} : $d_{pos}(a) = \max_{p \in \text{m\^eme cluster}} \|z_a - z_p\|_2$ --- force la compaction intra-cluster.
  \item \textbf{Semi-hard negative} : le n\'egatif le plus proche dont $d_{neg} > d_{pos}$. S'il n'existe pas de n\'egatif semi-hard, fallback sur le n\'egatif global le plus proche parmi les classes diff\'erentes.
  \item \textbf{S\'election des ancres} : seules les cellules avec $w_i \geq 1.2$ (poids de raret\'e suffisant) participent. Maximum 64 ancres par mini-batch.
\end{itemize}

\textbf{Formule compl\`ete de la Triplet Loss} :
\[
  \mathcal{L}_{triplet} = \frac{1}{|\mathcal{A}|} \sum_{a \in \mathcal{A}} \max\!\Bigl(0,\; \underbrace{\|z_a - z_p\|_2}_{d_{pos}(a)} - \underbrace{\|z_a - z_n\|_2}_{d_{neg}(a)} + m\Bigr)
\]
\textit{Description des termes :}
\begin{itemize}[nosep]
  \item $\mathcal{A}$ : ensemble des ancres retenues dans le mini-batch (filtr\'ees par \texttt{rare\_triplet\_min\_weight}).
  \item $z_a, z_p, z_n \in \mathbb{R}^{d_z}$ : embeddings de l'ancre, du positif et du n\'egatif.
  \item $d_{pos}(a)=\|z_a-z_p\|_2$ : distance ancre--positif.
  \item $d_{neg}(a)=\|z_a-z_n\|_2$ : distance ancre--n\'egatif.
  \item $m$ : marge de s\'eparation (\texttt{rare\_triplet\_margin}).
\end{itemize}
La loss est nulle quand $d_{neg} > d_{pos} + m$ (marge respect\'ee). Sinon, elle pousse \`a accro\^itre $d_{neg}$ ou \`a r\'eduire $d_{pos}$.

\textbf{Ramp-up $\rho(e)$ : mont\'ee progressive de la Triplet Loss.}

La Triplet Loss n'est pas activ\'ee brutalement \`a une \'epoque fixe : son poids effectif dans $\mathcal{L}_{total}$ est $\rho(e)\cdot\lambda_{rare}$. Dans le code, la dur\'ee de ramp-up n'est pas libre :
\[
  T_{\text{ramp}}=\max\!\bigl(1,\min(20,E_{\text{tot}}-e_{\text{start}})\bigr)
\]
et
\[
  \rho(e)=
  \begin{cases}
    0, & e<e_{\text{start}}\\
    \min\!\left(1,\dfrac{e-e_{\text{start}}}{T_{\text{ramp}}}\right), & e\ge e_{\text{start}}
  \end{cases}
\]
\textit{Description des termes :}
\begin{itemize}[nosep]
  \item $e$ : indice d'\'epoque courant.
  \item $e_{\text{start}}$ : \'epoque de d\'emarrage de la Triplet Loss (\texttt{rare\_triplet\_start\_epoch}).
  \item $E_{\text{tot}}$ : nombre total d'\'epoques d'entra\^inement (\texttt{epochs}).
  \item $T_{\text{ramp}}$ : dur\'ee effective du ramp-up, calcul\'ee comme ci-dessus.
  \item $\rho(e) \in [0,1]$ : coefficient multiplicatif appliqu\'e \`a $\lambda_{rare}$.
\end{itemize}

Cette fonction se d\'ecompose en trois r\'egimes distincts :

\begin{center}
\renewcommand{\arraystretch}{1.25}
\begin{tabular}{c c c c}
\toprule
\textbf{R\'egime} & \textbf{\'Epoques} & \textbf{$\rho(e)$} & \textbf{Poids effectif triplet} \\
\midrule
Inactif & $e < e_{\text{start}}$ & 0 & $0 \times \lambda_{rare} = \mathbf{0}$ (Triplet silencieuse) \\
Mont\'ee & $e_{\text{start}} \le e < e_{\text{start}} + T_{\text{ramp}}$ & $\dfrac{e-e_{\text{start}}}{T_{\text{ramp}}} \in\; ]0,\,1[$ & $\rho(e)\,\lambda_{rare}$ (croissance lin\'eaire) \\
Plateau & $e \ge e_{\text{start}} + T_{\text{ramp}}$ & 1 & $1 \times \lambda_{rare} = \lambda_{rare}$ (pleine force) \\
\bottomrule
\end{tabular}
\end{center}

\textbf{Valeurs cl\'es (forme g\'en\'erique)} : $\rho(e_{\text{start}})=0$, \quad $\rho\!\left(e_{\text{start}}+\tfrac{T_{\text{ramp}}}{2}\right)=0.5$, \quad $\rho(e_{\text{start}}+T_{\text{ramp}})=1$.

\begin{figure}[H]
\centering
\resizebox{0.9\linewidth}{!}{%
\begin{tikzpicture}[x=0.13cm, y=4cm]

% Phase backgrounds
\fill[phaseblue!10] (0, -0.06) rectangle (30, 1.12);
\fill[phaseorange!8] (30, -0.06) rectangle (120, 1.12);
\node[font=\small\bfseries, text=phaseblue, anchor=south] at (15, 1.12) {Warm-up (0--29)};
\node[font=\small\bfseries, text=phaseorange, anchor=south] at (75, 1.12) {Phase pond\'er\'ee (30--119)};

% Axes
\draw[-{Stealth}, thick] (-3, 0) -- (126, 0) node[right, font=\small] {\'Epoque $e$};
\draw[-{Stealth}, thick] (0, -0.06) -- (0, 1.25) node[above, font=\small] {$\rho(e)$};

% Y axis ticks
\foreach \y/\lab in {0/0, 0.25/0.25, 0.5/0.5, 0.75/0.75, 1.0/1} {
  \draw[gray!60] (-1.5, \y) -- (0, \y);
  \node[left, font=\scriptsize] at (-1.5, \y) {\lab};
  \draw[dotted, gray!30] (0, \y) -- (124, \y);
}

% X axis ticks (selected)
\foreach \x in {0,10,20,30,40,50,55,60,70,80,90,100,110,120} {
  \draw (\x, -0.02) -- (\x, 0.02);
  \node[below, font=\tiny] at (\x, -0.03) {\x};
}
\node[below, font=\tiny, text=tripletyellow!80!black] at (35, -0.06) {\textbf{35}};

% ρ(e) curve
% Warmup: inactive (dashed gray)
\draw[gray!50, dashed, thick] (0, 0) -- (30, 0);
% Phase 2 before triplet start: ρ=0
\draw[tripletyellow!80!black, very thick] (30, 0) -- (35, 0);
% Ramp-up: linear 35→55
\draw[tripletyellow!80!black, very thick] (35, 0) -- (55, 1);
% Plateau: ρ=1, 55→120
\draw[tripletyellow!80!black, very thick] (55, 1) -- (120, 1);

% Hatched ramp-up zone
\fill[tripletyellow!12] (35, 0) -- (55, 1) -- (55, 0) -- cycle;

% Key point dots
\fill[tripletyellow!80!black] (35,  0   ) circle (2.5pt);
\fill[tripletyellow!80!black] (40,  0.25) circle (2.5pt);
\fill[tripletyellow!80!black] (45,  0.5 ) circle (2.5pt);
\fill[tripletyellow!80!black] (50,  0.75) circle (2.5pt);
\fill[tripletyellow!80!black] (55,  1.0 ) circle (2.5pt);

% Annotations
\node[above right, font=\tiny, text=tripletyellow!80!black] at (35, 0.02)
  {$\rho{=}0$\;$\lambda_{eff}{=}0$};
\node[right, font=\tiny, text=tripletyellow!80!black] at (46, 0.5)
  {$\rho{=}0.5$\;$\lambda_{eff}{=}\lambda_{rare}/2$};
\node[above, font=\tiny, text=tripletyellow!80!black] at (55, 1.01)
  {$\rho{=}1$\;$\lambda_{eff}{=}\lambda_{rare}$};

% Dashed verticals at key epochs
\draw[dashed, tripletyellow!60] (35, 0) -- (35, -0.05);
\draw[dashed, tripletyellow!60] (55, 0) -- (55, -0.05);

% Effective lambda on right side
\draw[gray!60, -{Stealth}] (122, 0) -- (122, 1.1) node[right, font=\tiny, gray] {$\lambda_{eff}$};
\node[right, font=\tiny, gray] at (122, 0)   {0};
\node[right, font=\tiny, gray] at (122, 1.0) {$\lambda_{rare}$};

% Annotation: ramp zone
\draw[{|[scale=1.2]}-{|[scale=1.2]}, thick, tripletyellow!60] (35, -0.05) -- (55, -0.05)
  node[midway, below, font=\tiny, text=tripletyellow!80!black] {$T_{\text{ramp}}$ \'ep. de mont\'ee};

\end{tikzpicture}%
}
\caption{Profil du coefficient de ramp-up $\rho(e)$ (illustration avec valeurs par d\'efaut : $e_{\text{start}}=35$, \texttt{epochs}=120, donc $T_{\text{ramp}}=20$). Le poids effectif de la Triplet Loss dans $\mathcal{L}_{total}$ est $\rho(e)\cdot\lambda_{rare}$ : nul avant l'\'epoque de d\'emarrage, en mont\'ee lin\'eaire sur la p\'eriode de ramp-up, puis stable \`a $\lambda_{rare}$. La zone hachet\'ee repr\'esente la p\'eriode de transition.}
\label{fig:rampup}
\end{figure}

\textbf{Pourquoi un ramp-up et non une activation brutale ?} Si on activait la Triplet Loss \`a pleine force ($\lambda_{rare}$) d\`es $e_{\text{start}}$, le gradient triplet perturberait brusquement une reconstruction NB d\'ej\`a en cours de convergence. Les triplets form\'es \`a partir de pseudo-labels encore imparfaits g\'en\`ereraient des signaux contradictoires, pouvant d\'estabiliser l'espace latent. Le ramp-up laisse au mod\`ele le temps de s'adapter progressivement au nouveau signal.

\textbf{Pourquoi d\'emarrer apr\`es le warm-up ?} Quelques \'epoques suppl\'ementaires de pond\'eration pure entre la transition et $e_{\text{start}}$ permettent aux pseudo-labels de se stabiliser avant l'activation de la Triplet Loss. D\'emarrer trop t\^ot (pseudo-labels instables) ou trop tard (structure latente d\'ej\`a fig\'ee) peut nuire aux performances.

\begin{parambox}{Param\`etres Triplet Loss --- valeurs par d\'efaut}
\begin{tabular}{lll}
\texttt{rare\_loss\_type} & \textbf{triplet} & Type de r\'egularisation rare \\
\texttt{rare\_triplet\_weight} ($\lambda_{rare}$) & \textbf{0.1} & Poids dans la perte totale \\
\texttt{rare\_triplet\_margin} ($m$) & \textbf{0.4} & Marge de s\'eparation \\
\texttt{rare\_triplet\_min\_weight} & \textbf{1.2} & Seuil de poids pour \^etre ancre \\
\texttt{rare\_triplet\_start\_epoch} & \textbf{35} & \'Epoque d'activation absolue \\
\texttt{ramp-up} & \textbf{20 (par d\'efaut)} & Dur\'ee effective : $\max(1,\min(20,\texttt{epochs} - \texttt{rare\_triplet\_start\_epoch}))$ \\
\texttt{max\_triplet\_anchors\_per\_batch} & \textbf{64} & Ancres max par mini-batch \\
\end{tabular}
\end{parambox}

% =============================================================================
\section{Phase de clustering}
\label{sec:clustering}
% =============================================================================

Les embeddings latents $z \in \mathbb{R}^{128}$ sont extraits en fin d'entra\^inement (\texttt{cluster\_preprocess\_mode=none} : pas de r\'eduction dimensionnelle suppl\'ementaire). HDBSCAN est appliqu\'e directement sur ces 128 dimensions.

HDBSCAN effectue ensuite l'assignation des clusters en estimant la densit\'e locale hi\'erarchique.

\begin{figure}[H]
\centering
\resizebox{\linewidth}{!}{%
\begin{tikzpicture}[
  node distance=0.8cm and 1.5cm,
  step/.style={rectangle, draw, thick, rounded corners=5pt, align=center,
    minimum width=3cm, minimum height=1.2cm, font=\footnotesize},
  arr/.style={-{Stealth[length=2.5mm]}, thick},
  param/.style={font=\tiny\ttfamily, text=gray!70},
]

\node[step, fill=latentcolor!15] (embed) {\textbf{Embeddings latents}\\$z \in \mathbb{R}^{n \times 128}$\\(sortie de l'encoder)};

\node[step, fill=clustergreen!10, right=2cm of embed, text width=3.5cm] (hdbscan) {\textbf{HDBSCAN}\\S\'election EOM\\Densit\'e hi\'erarchique};

\node[step, fill=clustergreen!20, right=2cm of hdbscan, text width=3cm] (reassign) {\textbf{R\'eassignation}\\Bruit $\to$ plus proche\\centroid};

\node[right=1.5cm of reassign, font=\small\bfseries, text=clustergreen, align=center] (out) {$k$ clusters\\d\'etect\'es};

% Arrows
\draw[arr] (embed) -- node[above, font=\tiny]{Direct (pas de PCA)} (hdbscan);
\draw[arr] (hdbscan) -- node[above, font=\tiny]{Bruit = $-1$} (reassign);
\draw[arr, clustergreen] (reassign) -- (out);

% Parameters
\node[param, below=0.4cm of hdbscan, text width=3.5cm, align=center] {%
  \texttt{min\_cluster\_size = 4}\\
  \texttt{min\_samples = 2}\\
  \texttt{selection = eom}};

\end{tikzpicture}%
}
\caption{Pipeline de clustering. HDBSCAN est appliqu\'e directement sur les embeddings 128D, sans r\'eduction dimensionnelle. Les points class\'es \og bruit \fg{} ($-1$) sont r\'eassign\'es au centroid de cluster le plus proche.}
\label{fig:clustering}
\end{figure}

\textbf{Justification des param\`etres \texttt{min\_cluster\_size = 4} et \texttt{min\_samples = 2}}

\begin{itemize}[nosep]
  \item \texttt{min\_cluster\_size = 4} : permet de d\'etecter des clusters aussi petits que 4 cellules, ce qui est important pour les populations rares.
  \item \texttt{min\_samples = 2} : abaisse le seuil de \og point core \fg{}, r\'eduisant le nombre de cellules class\'ees comme bruit. Avec \texttt{min\_samples = 3} (d\'efaut HDBSCAN), les micro-clusters seraient enti\`erement class\'es bruit.
  \item \texttt{selection = eom} (Excess of Mass) : pr\'eserve les petits clusters \`a c\^ot\'e des grands, contrairement \`a \texttt{leaf} qui fragmenterait davantage.
\end{itemize}

\begin{parambox}{Param\`etres HDBSCAN --- valeurs par d\'efaut}
\begin{tabular}{lll}
\texttt{clustering\_method} & \textbf{hdbscan} & M\'ethode de clustering final \\
\texttt{hdbscan\_min\_cluster\_size} & \textbf{4} & Taille minimale d'un cluster \\
\texttt{hdbscan\_min\_samples} & \textbf{2} & Points cores minimum \\
\texttt{hdbscan\_cluster\_selection\_method} & \textbf{eom} & Excess of Mass \\
\texttt{hdbscan\_reassign\_noise} & \textbf{true} & R\'eassigner le bruit \\
\texttt{cluster\_preprocess\_mode} & \textbf{none} & Pas de post-traitement \\
\end{tabular}
\end{parambox}

% =============================================================================
\section{Diagramme r\'ecapitulatif des flux de gradients}
\label{sec:gradient_flow}
% =============================================================================

Ce diagramme montre comment les diff\'erents signaux de gradient influencent l'encoder pendant la phase 2.

\textbf{Signal 1 --- Reconstruction NB pond\'er\'ee} (signal principal) : le d\'ecodeur tente de reconstruire les raw counts $x_{ij}$ depuis $z$. Les gradients $\nabla_z \mathcal{L}_{NB}$ atteignent l'encoder via backpropagation. Gr\^ace aux poids $w_i$, les cellules rares g\'en\`erent une contribution plus forte au gradient, poussant l'encoder \`a mieux discriminer leur profil d'expression.

\textbf{Signal 2 --- Triplet Loss} ($\lambda_{rare}$, \`a partir de l'\'ep.~$e_{\text{start}}$) : les gradients $\nabla_z \mathcal{L}_{triplet}$ s'appliquent directement sur $z$ (pas sur le d\'ecodeur). Ils imposent une \textbf{contrainte de s\'eparation locale} : les cellules rares (ancres) doivent \^etre plus proches de leurs alli\'ees que de leurs intrus dans $\mathbb{R}^{128}$, ind\'ependamment de la qualit\'e de reconstruction.

\textbf{Interaction entre les deux signaux} : la reconstruction pond\'er\'ee agit globalement (tous les g\`enes de toutes les cellules), tandis que la Triplet Loss est locale (quelques triplets de cellules rares par mini-batch). Le param\`etre $\lambda_{rare}$ contr\^ole le rapport de force entre ces deux signaux : trop fort, la Triplet Loss d\'estabilise la reconstruction ; trop faible, elle ne structure pas assez l'espace latent. En configuration multi-batch, le gradient DANN s'ajoute comme troisi\`eme signal.

\begin{figure}[H]
\centering
\resizebox{\linewidth}{!}{%
\begin{tikzpicture}[
  node distance=1.0cm and 1.5cm,
  module/.style={rectangle, draw, thick, rounded corners=6pt, align=center,
    minimum width=2.8cm, minimum height=1.1cm, font=\footnotesize},
  loss/.style={ellipse, draw, thick, align=center, font=\footnotesize, inner sep=3pt},
  arr/.style={-{Stealth[length=2.5mm]}, thick},
  garr/.style={-{Stealth[length=2.5mm]}, thick, dashed},
]

% Forward path
\node[module, fill=enccolor!12] (enc) {\textbf{Encodeur}\\512$\to$256$\to$128};
\node[module, fill=latentcolor!15, right=1.5cm of enc] (z) {\textbf{$z$}\\128 dim};
\node[module, fill=deccolor!12, right=1.5cm of z] (dec) {\textbf{D\'ecodeur}\\128$\to$256$\to$512};

% Losses
\node[loss, fill=nbblue!12, right=1.5cm of dec] (lnb) {$\mathcal{L}_{NB}$\\pond\'er\'ee};
\node[loss, fill=tripletyellow!15, below=1.2cm of z] (ltri) {$\mathcal{L}_{triplet}$\\$\lambda=\lambda_{rare}$};

% Forward arrows
\draw[arr] (enc) -- (z);
\draw[arr] (z) -- (dec);
\draw[arr] (dec) -- (lnb);

% Gradient arrows (backward)
\draw[garr, nbblue, line width=1.5pt] (lnb) -- node[above, font=\tiny, text=nbblue]{$\nabla$ recon} (dec);
\draw[garr, nbblue, line width=1.5pt] (dec) -- (z);
\draw[garr, nbblue, line width=1.5pt] (z) -- (enc);

% Triplet gradient
\draw[arr] (z) -- (ltri);
\draw[garr, tripletyellow!80!black, line width=1.5pt] (ltri) -- node[right, font=\tiny, text=tripletyellow!80!black, text width=1.5cm]{$\nabla$ rare\\(semi-hard)} (z);

% Legend
\node[draw, rounded corners=4pt, fill=white, font=\tiny, align=left, text width=5cm]
  at ($(enc)+(0,-3.5)$) {%
  \textbf{L\'egende des gradients :}\\
  \textcolor{nbblue}{\rule{0.5cm}{2pt}} Reconstruction NB (principal)\\
  \textcolor{tripletyellow!80!black}{\rule{0.5cm}{2pt}} Triplet Loss (structuration locale)\\
  \textcolor{dannred}{\rule{0.5cm}{2pt}} DANN (optionnel, multi-batch)};

\end{tikzpicture}%
}
\caption{Flux de gradients pendant la phase 2. Deux signaux principaux convergent vers l'encoder : (1) la reconstruction NB pond\'er\'ee (signal principal), (2) la Triplet Loss qui structure les cellules rares dans $z$. En multi-batch, le DANN ajoute un troisi\`eme signal.}
\label{fig:gradients}
\end{figure}

% =============================================================================
\section{R\'ecapitulatif complet des param\`etres}
\label{sec:all_params}
% =============================================================================

\begin{longtable}{p{4.8cm}p{2.2cm}p{7cm}}
\caption{Param\`etres de scRAW et leurs valeurs par d\'efaut} \label{tab:all_params} \\
\toprule
\textbf{Param\`etre} & \textbf{Valeur} & \textbf{Description} \\
\midrule
\endfirsthead
\toprule
\textbf{Param\`etre} & \textbf{Valeur} & \textbf{Description} \\
\midrule
\endhead
\midrule
\multicolumn{3}{r}{\textit{Suite page suivante\ldots}} \\
\bottomrule
\endfoot
\bottomrule
\endlastfoot

\multicolumn{3}{l}{\cellcolor{enccolor!8}\textbf{Architecture r\'eseau}} \\
\midrule
\texttt{z\_dim} & 128 & Dimension de l'espace latent \\
\texttt{encoder\_layers} & 512,256,128 & Couches cach\'ees de l'encoder \\
\texttt{decoder\_layers} & (sym\'etrique) & Miroir de l'encoder \\
\texttt{activation} & leaky\_relu & Fonction d'activation \\
\texttt{dropout} & 0.1 & Taux de dropout \\

\midrule
\multicolumn{3}{l}{\cellcolor{phaseblue!8}\textbf{Entra\^inement}} \\
\midrule
\texttt{epochs} & 120 & \'Epoques totales \\
\texttt{warmup\_epochs} & 30 & \'Epoques de warm-up \\
\texttt{batch\_size} & 256 & Taille des mini-batches \\
\texttt{lr} & 0.001 & Taux d'apprentissage (Adam) \\
\texttt{scheduler} & CosineAnnealingLR & $\eta_{min} = \texttt{lr} \times 0.01$ \\
\texttt{gradient clipping} & max\_norm=5.0 & Clip des gradients par norme \\
\texttt{random\_state} & 42 & Graine al\'eatoire \\

\midrule
\multicolumn{3}{l}{\cellcolor{nbblue!8}\textbf{Reconstruction}} \\
\midrule
\texttt{reconstruction\_distribution} & nb & Distribution Negative Binomial \\
\texttt{nb\_theta} & 10.0 & Dispersion NB (fixe) \\
\texttt{nb\_input\_transform} & log1p & Transformation d'entr\'ee encoder \\
\texttt{masking\_rate} & 0.2 & Masquage 20\% (warm-up seul) \\
\texttt{masked\_recon\_weight} & 0.75 & Poids r\'eduit positions masqu\'ees \\
\texttt{masking\_apply\_weighted} & false & Masquage d\'esactiv\'e en phase 2 \\

\midrule
\multicolumn{3}{l}{\cellcolor{weightpurple!8}\textbf{Pond\'eration dynamique}} \\
\midrule
\texttt{pseudo\_label\_method} & leiden & M\'ethode de pseudo-labeling \\
\texttt{weight\_exponent} ($\alpha$) & 0.4 & Exposant inverse-fr\'equence \\
\texttt{density\_knn\_k} & 15 & Voisins pour densit\'e locale \\
\texttt{density\_weight\_exponent} ($\beta$) & 1.0 & Exposant distance kNN \\
\texttt{density\_weight\_clip} & 5.0 & Clip max densit\'e \\
\texttt{weight\_fusion\_mode} & additive & Mode de fusion \\
\texttt{cluster\_density\_alpha} ($\gamma$) & 0.6 & Ratio fr\'equence/densit\'e \\
\texttt{min\_cell\_weight} & 0.25 & Poids minimum \\
\texttt{max\_cell\_weight} & 10.0 & Poids maximum \\
\texttt{dynamic\_weight\_update\_interval} & 10 & Recalcul tous les $N$ cycles \\
\texttt{dynamic\_weight\_momentum} ($\mu$) & 0.7 & Inertie EMA \\

\midrule
\multicolumn{3}{l}{\cellcolor{tripletyellow!12}\textbf{R\'egularisation rare (Triplet)}} \\
\midrule
\texttt{rare\_loss\_type} & triplet & Type de perte rare \\
\texttt{rare\_triplet\_weight} ($\lambda_{rare}$) & 0.1 & Poids dans la perte totale \\
\texttt{rare\_triplet\_margin} ($m$) & 0.4 & Marge de s\'eparation \\
\texttt{rare\_triplet\_min\_weight} & 1.2 & Seuil ancre \\
\texttt{rare\_triplet\_start\_epoch} & 35 & \'Epoque d'activation \\
\texttt{max\_triplet\_anchors\_per\_batch} & 64 & Ancres max par batch \\

\midrule
\multicolumn{3}{l}{\cellcolor{clustergreen!8}\textbf{Clustering HDBSCAN}} \\
\midrule
\texttt{clustering\_method} & hdbscan & M\'ethode de clustering final \\
\texttt{hdbscan\_min\_cluster\_size} & 4 & Taille minimale de cluster \\
\texttt{hdbscan\_min\_samples} & 2 & Points cores minimum \\
\texttt{hdbscan\_cluster\_selection\_method} & eom & Excess of Mass \\
\texttt{hdbscan\_reassign\_noise} & true & R\'eassigner le bruit \\
\texttt{cluster\_preprocess\_mode} & none & Pas de post-traitement \\

\midrule
\multicolumn{3}{l}{\cellcolor{dannred!8}\textbf{Correction de batch (optionnelle)}} \\
\midrule
\texttt{use\_batch\_conditioning} & false & Conditionner le d\'ecodeur sur le batch \\
\texttt{adversarial\_batch\_weight} ($\lambda_{adv}$) & 0.0 & Poids DANN (0 = d\'esactiv\'e) \\
\texttt{adversarial\_lambda} & 1.0 & \'Echelle du gradient reversal \\
\texttt{adversarial\_start\_epoch} & 0 & \'Epoque de d\'emarrage du DANN \\
\texttt{adversarial\_ramp\_epochs} & 0 & Dur\'ee du ramp-up DANN \\
\texttt{mmd\_batch\_weight} & 0.0 & Poids MMD (0 = d\'esactiv\'e) \\


\end{longtable}

\begin{notebox}
Le tableau ci-dessus pr\'esente les \textbf{valeurs par d\'efaut} des param\`etres de scRAW. Ces valeurs peuvent \^etre ajust\'ees selon le jeu de donn\'ees. scRAW attend des donn\'ees pr\'etrait\'ees en entr\'ee (normalisation, log1p, HVG, mise \`a l'\'echelle).
\end{notebox}



% =============================================================================
\section{Diagramme r\'ecapitulatif complet}
\label{sec:full_diagram}
% =============================================================================

\begin{figure}[H]
\centering
\resizebox{\linewidth}{!}{%
\begin{tikzpicture}[
  node distance=1.0cm and 0.8cm,
  stage/.style={rectangle, draw, thick, rounded corners=6pt, align=center,
    minimum width=4cm, minimum height=1.4cm, font=\small},
  detail/.style={rectangle, draw=gray!40, rounded corners=3pt, fill=gray!5,
    align=left, font=\tiny, text width=4.5cm, inner sep=3pt},
  arr/.style={-{Stealth[length=3mm]}, thick, gray!70},
  darr/.style={-{Stealth[length=3mm]}, thick, dashed},
]

% Main pipeline (vertical)
\node[stage, fill=gray!8] (input) {\textbf{Donn\'ees scRNA-seq}\\$n$ cellules $\times$ $G$ g\`enes};

\node[stage, fill=blue!6, below=of input] (preproc) {\textbf{Pr\'etraitement (Scanpy)}\\Norm + HVG + Scale};

\node[stage, fill=phaseblue!10, below=of preproc] (warmup) {\textbf{Phase 1 : Warm-up}\\30 \'epoques (0--29)};

\node[stage, fill=phaseorange!12, below=of warmup] (weighted) {\textbf{Phase 2 : Pond\'er\'ee}\\90 \'epoques (30--119)};

\node[stage, fill=clustergreen!12, below=of weighted] (cluster) {\textbf{Phase 3 : HDBSCAN}\\sur espace latent 128D};

\node[stage, fill=successgreen!15, below=of cluster] (output) {\textbf{Labels finaux}\\$k$ clusters d\'etect\'es};

% Main arrows
\draw[arr] (input) -- (preproc);
\draw[arr] (preproc) -- (warmup);
\draw[arr] (warmup) -- (weighted);
\draw[arr] (weighted) -- (cluster);
\draw[arr] (cluster) -- (output);

% Detail boxes
\node[detail, right=1.5cm of warmup] (warmup_d) {%
  $\bullet$ NB uniforme ($w_i=1$), 30 \'ep.\\
  $\bullet$ Masquage 20\% actif (per mini-batch)\\
  $\bullet$ DANN optionnel (multi-batch)\\
  $\bullet$ lr=0.001, batch=256};

\node[detail, right=1.5cm of weighted] (weighted_d) {%
  $\bullet$ NB pond\'er\'ee : $\sum w_i \mathcal{L}_{NB}$\\
  $\bullet$ Poids dynamiques $\Delta$10 \'ep. EMA($\mu$)\\
  $\bullet$ Triplet d\`es \'ep. $e_{\text{start}}$ ($\lambda_{rare}$)\\
  $\bullet$ Masquage d\'esactiv\'e en phase 2\\
  $\bullet$ DANN optionnel (multi-batch)};

\node[detail, right=1.5cm of cluster] (cluster_d) {%
  $\bullet$ HDBSCAN c4s2, EOM\\
  $\bullet$ Pas de post-traitement\\
  $\bullet$ Bruit $\to$ plus proche cluster\\
  $\bullet$ $k$ d\'etect\'e automatiquement};

% Arrows to details
\draw[gray!40, thick] (warmup.east) -- (warmup_d.west);
\draw[gray!40, thick] (weighted.east) -- (weighted_d.west);
\draw[gray!40, thick] (cluster.east) -- (cluster_d.west);

% Weight feedback
\node[detail, left=1.5cm of weighted, fill=weightpurple!5, draw=weightpurple!30, text width=3.8cm] (weight_d) {%
  \textbf{Pond\'eration bicrit\`ere}\\
  $\bullet$ Leiden $\to$ pseudo-labels\\
  $\bullet$ $w_{freq} = (1/f)^{\alpha}$\\
  $\bullet$ $w_{dens} = (d_{kNN}/\text{m\'ed})^{\beta}$\\
  $\bullet$ Fusion : $\gamma \cdot freq + (1-\gamma) \cdot dens$\\
  $\bullet$ Momentum $\mu$, clip [0.25, 10]};

\draw[darr, weightpurple] (weight_d.east) -- (weighted.west);
\draw[darr, weightpurple] (weighted.south west) |- ++(0,-0.3) -| (weight_d.south) node[near start, below, font=\tiny, text=weightpurple]{MAJ $\Delta$10 \'ep.};

\end{tikzpicture}%
}
\caption{Diagramme r\'ecapitulatif complet du pipeline scRAW avec tous les param\`etres cl\'es.}
\label{fig:full_diagram}
\end{figure}

% =============================================================================
\section{R\'esum\'e des caract\'eristiques}
\label{sec:conclusion}
% =============================================================================

scRAW repose sur les composants suivants :

\begin{enumerate}[nosep]
  \item \textbf{Reconstruction} : distribution Negative Binomial sur les raw counts, avec entr\'ee transform\'ee (log1p) et masquage stochastique pendant la phase de warm-up.
  \item \textbf{Pond\'eration dynamique} : calcul bicrit\`ere combinant fr\'equence de pseudo-cluster et densit\'e locale, mis \`a jour p\'eriodiquement avec stabilisation EMA.
  \item \textbf{Contrainte latente} : Triplet Loss avec semi-hard mining, ciblant les cellules rares, activ\'ee progressivement via un ramp-up lin\'eaire.
  \item \textbf{Correction batch} : module DANN optionnel, activable en configuration multi-batch.
  \item \textbf{Clustering} : HDBSCAN appliqu\'e directement sur l'espace latent, avec r\'eassignation du bruit.
\end{enumerate}

L'ensemble de ces m\'ecanismes permet \`a scRAW de produire des repr\'esentations latentes adapt\'ees au clustering de donn\'ees scRNA-seq, avec une attention particuli\`ere \`a la d\'etection des populations cellulaires rares.

\end{document}
